\section{Nghiên cứu liên quan}\label{sec:related}
% Trước khi viết bản tiếng Anh: mở PDF trong ../refs/ để đối chiếu từng câu trích.

\paragraph{Tấn công và benchmark.} Greshake và cộng sự~\cite{greshake2023indirect} đưa ra khái niệm chèn lệnh gián
tiếp; Liu và cộng sự~\cite{liu2024formalizing} hình thức hóa tấn công và phòng thủ chèn lệnh. BIPIA~\cite{yi2025bipia}
đo chèn lệnh gián tiếp vào LLM; InjecAgent~\cite{zhan2024injecagent}, AgentDojo~\cite{debenedetti2024agentdojo} và
ASB~\cite{zhang2025asb} chuyển sang tác tử dùng công cụ, nơi tấn công gây ra hành động. Các bộ này đều bằng tiếng
Anh; Deng và cộng sự~\cite{deng2024multijail} cho thấy cơ chế an toàn của LLM yếu hơn ở ngôn ngữ ít tài nguyên,
nhưng trong bối cảnh jailbreak chứ không phải tác tử.

\paragraph{Phòng thủ.} Ở mức prompt có spotlighting~\cite{hines2024spotlighting} và lời nhắc kiểu sandwich.
StruQ~\cite{chen2025struq} và SecAlign~\cite{chen2025secalign} huấn luyện lại mô hình để tách chỉ thị khỏi dữ liệu.
Bộ phân loại như Prompt Guard~2~\cite{meta2025promptguard2,chennabasappa2025llamafirewall} sàng lọc đầu vào.
MELON~\cite{zhu2025melon} so hành động của lượt chạy gốc với lượt chạy lại có che yêu cầu; Progent~\cite{shi2025progent}
áp chính sách đặc quyền lên lời gọi công cụ. Hướng an toàn theo thiết kế gồm CaMeL~\cite{debenedetti2026camel},
kiểm soát luồng thông tin FIDES~\cite{costa2025fides} và các mẫu thiết kế trong~\cite{beurerkellner2025patterns}.
\ours{} gần hướng này nhất: giống CaMeL, nó không sửa mô hình, nhưng thay vì tách luồng điều khiển và chạy chương
trình trong một trình thông dịch riêng, \ours{} chỉ áp luật chuỗi lên quỹ đạo nên gắn được vào tác tử ReAct sẵn có. Tấn công thích nghi vượt qua nhiều phòng thủ đã
công bố~\cite{zhan2025adaptive,nasr2025attacker}; chúng tôi báo cáo riêng một tập tấn công thích nghi
(Mục~\ref{sec:external}).
